\documentclass[11pt,a4paper]{article}
\usepackage[margin=25mm]{geometry}
\usepackage{amsmath,amssymb,amsthm,booktabs}
\newtheorem{proposition}{Proposition}
\title{Two oriented square-torus amalgamations}
\author{Human Compression project}\date{}
\begin{document}\maketitle
The fixed Schulte--Weiss survey attributes the following two examples to
Hubard--Weiss, \emph{Self-duality of chiral polytopes} (2005), p.~132.
We verify their presented groups and the actual intersection and chirality
conditions. This is a recovery of two established examples, not a new
existence result or a classification of toroidal amalgamations.

Write $T(s,t)=\{4,4\}_{(s,t)}$ with its specified orientation, and put
\[
 R(x,y;s,t)=(x^{-1}y)^s(xy^{-1})^t.
\]
We retain two precise standard interfaces stated in the survey: the
oriented square-torus presentation on rotations $x,y$ is
$x^4=y^4=(xy)^2=R(x,y;s,t)=1$; and the rank-four rotation reconstruction
criterion consists of the string rotation relations and the three
intersection equalities below. The square-torus convention is the lattice
spanned by $(s,t)$ and $(-t,s)$. In the cases here its index is ten, so
its rotation group has order forty. No completeness assertion for other
parameters is used.

\begin{proposition}
For $(s,t)=(1,3)$ or $(3,1)$, let
\begin{align*}
 G(s,t)=\langle u,v,w\mid{}&u^4=v^4=w^4=(uv)^2=(vw)^2=(uvw)^2=1,\\
                         &R(u,v;s,t)=R(v,w;1,3)=1\rangle .
\end{align*}
These are the automorphism groups of the oriented universal chiral
amalgamations with facets $T(s,t)$ and vertex figures $T(1,3)$.
Their orders are respectively $2000$ and $960$; hence the two resulting
polytopes are not isomorphic.
\end{proposition}
\begin{proof}
The supplied finite calculation constructs an isomorphism from each full
finitely presented group to a permutation group, of degrees fifty and
twenty respectively. This is not merely a finite quotient representation.
The following quantities are checked in those actual permutation groups:
\[
\begin{array}{c|r|r|r|r}
(s,t)&|G(s,t)|&|\langle u,v\rangle|&|\langle v,w\rangle|
 &|\langle u,v\rangle\cap\langle v,w\rangle|\\\hline
(1,3)&2000&40&40&4\\
(3,1)&960&40&40&4
\end{array}
\]
Each of $u,v,w$ has order four, and
\[
 \langle u\rangle\cap\langle v\rangle
 =\langle v\rangle\cap\langle w\rangle=1,\qquad
 \langle u,v\rangle\cap\langle v,w\rangle=\langle v\rangle.
\]
The local finitely presented groups are separately checked to have order
forty. Their natural epimorphisms onto the actual end subgroups are
therefore isomorphisms, so the prescribed toroids do not collapse.
The retained reconstruction criterion now supplies a rank-four polytope
with precisely those oriented local types.

The regularity test would require an automorphism taking
\[
              (u,v,w)\longmapsto(u^{-1},u^2v,w).
\]
After this substitution, both displayed toroidal relators are nonidentity
in each group. Such an automorphism cannot exist. The reconstructed
polytopes are therefore chiral and $G(s,t)$ is their full automorphism
group, rather than only the rotation subgroup of a regular polytope.

The two added relators are exactly the local defining relations beyond
the standard rotation relations. Every oriented amalgamation with these
specified facets and vertex figures consequently has its marked rotation
group as a quotient of $G(s,t)$. Together with the verified existence
and uncollapsed local types, the retained oriented amalgamation interface
gives the stated universal object. Finally unequal automorphism-group
orders exclude an isomorphism between the two polytopes.
\end{proof}

\paragraph{Verification and attribution.}
The executable calculation uses GAP's isomorphism algorithm for finite
presentations. It checks bijectivity and all defining relations, and
exports the three literal permutations. A separate Python calculation,
using only composition and exhaustive closure of the resulting small
groups, verifies both group orders, every stated subgroup intersection,
local orders and both failed reflection relators. This second calculation
does not itself prove the upper order of the finite presentation; that
step remains the explicit GAP computation. A damaged marking is rejected.
These are computational checks, not Lean certificates.

The source is the fixed arXiv edition math/0608397v1, Section~4 (physical
p.~7) for the orders and attribution, and Section~5 (physical p.~12)
for the literal toroidal relator. The rank-four reconstruction and
reflection criteria occur in Section~2 (physical p.~4), attributed there
to Schulte--Weiss, \emph{Chirality and projective linear groups} (1994).
The cited 1994 and 2005 primary papers were not independently inspected
for this component; the exact interfaces above are retained from the
read survey, while the finite group and intersection calculations are
reproduced here. General chiral amalgamation existence or finiteness
questions are not resolved by these two examples.
\end{document}
